Nonlinear finite-element models are often defined by compact energy or
constitutive programs, while their solution requires distributed sparse
matrices, Krylov methods, and nonlinear globalization. Automatic
differentiation (AD) removes many manual derivative formulas, but it does not
remove the choice of where differentiation occurs. A tangent may be obtained
from the complete element energy, from a low-dimensional constitutive density,
or from colored actions of the assembled Hessian. Each route exposes a
different combination of local arithmetic, memory traffic, sparse insertion,
and communication.

This choice is scientific only if the compared routes represent the same
discrete problem. Quadrature, essential constraints, affine lifts, active
constitutive branches, engineering-shear conventions, and state ordering must
remain fixed. Distributed colored recovery introduces further requirements:
the structural pattern must be complete, same-color columns must not interfere
in an owned row, and every ghost value needed by a local action must be
available. Agreement of endpoint energies alone does not verify these
conditions.

Existing frameworks establish the principal component technologies. High-level
finite-element systems automate variational forms and parallel assembly
\citep{logg2012fenicsbook,baratta2025dolfinx}; JAX-based methods provide
differentiable mechanics and inverse problems
\citep{xue2023jaxfem,xue2026implicit}; and bridge architectures combine
differentiable local programs with established sparse solvers
\citep{yashchuk2023bringing,latyshev2025externaloperators,cattaneo2026jetsci}.
Recent work also develops locality-aware element differentiation, global AD
with sparse tangent materialization, and automatic sparsity detection
\citep{mahmoud2025locality,pundir2026tatva,hill2025sparser}. Consequently, the
novelty of the present study is not JAX--PETSc coupling, constitutive AD, or
graph coloring in isolation.

We instead ask two narrower questions. First, which assumptions guarantee that
element AD, constitutive AD, and colored sparse recovery differentiate one
quadrature-defined potential? Second, which fixed-state and distributed tests
expose errors in branch selection, local contraction, sparse recovery, or
Message Passing Interface (MPI) ownership? Both questions admit mathematical
statements and finite correctness tests. Performance selection is outside the
scope of the present study.

\subsection{Contributions}
\label{subsec:intro-contributions}

The paper makes four contributions.
\begin{enumerate}
  \item We formulate a branchwise equivalence result for element and
        constitutive differentiation of a fixed quadrature potential. The
        statement includes the constrained affine map and excludes branch
        switches and unresolved repeated spectra.
  \item We state the interference, overlap, and row-ownership conditions under
        which colored Hessian actions recover the same distributed sparse
        operator.
  \item We provide an evidence hierarchy that separates analytic and
        manufactured verification, canonical fixed-state MPI checks, and
        bounded nonlinear-outcome studies. Each numerical statement is tied to
        the algebraic object and admission conditions actually tested.
  \item We exercise this hierarchy on smooth scalar energies,
        orientation-preserving hyperelasticity, and a synthetic
        branch-structured Mohr--Coulomb endpoint potential. The latter is used
        to test derivative placement and is not presented as incremental
        plasticity validation.
\end{enumerate}

The resulting evidence is deliberately asymmetric. The manufactured and
analytic verifiers pass for the tested smooth problems. Element and
constitutive differentiation agree on 15 prescribed finite-element states, and
distributed colored recovery matches the prescribed element-route tangent
actions for $P_1(L_1)$ and $P_2(L_1)$ on one, two, and four ranks. Direct
complete sparse-matrix comparisons additionally pass at one rank. A separate
checksum-verified workstation study corroborates fixed-state tangent-action
agreement. We report no route timing and claim neither a fastest route nor a
crossover or scaling law.

Section~\ref{sec:related} positions this scope against the closest software and
methodological work. Sections~\ref{sec:methodology} and
\ref{sec:implementation} state the mathematical and distributed contracts.
Section~\ref{sec:benchmarks} defines the retained problems, and
Section~\ref{sec:validation} presents independent verification designs and
results. The remaining numerical outcomes and their limitations appear in
Section~\ref{sec:performance-solver-behavior}; Sections~\ref{sec:discussion}
and~\ref{sec:conclusion} summarize the resulting claim boundary.
